\documentclass[12pt]{article}

\input{../../../_assets/preambulo}

\begin{document}

    % 1. Foto de fondo
    % 2. Título
    % 3. Encabezado Izquierdo
    % 4. Color de fondo
    % 5. Coord x del titulo
    % 6. Coord y del titulo
    % 7. Fecha

    
	\input{../../../_assets/portada}
	\portadaExamenFotoDif{../../sapienza.jpg}{Computer Vision\\Examen V}{Computer Vision. Examen V}{MidnightBlue}{-8}{28}{2025-2026}{Joaquín Avilés de la Fuente}

    \begin{description}
        \item[Asignatura] Computer Vision.
        \item[Curso Académico] 2025-2026.
        \item[Grado] Grado en Informática.
        \item[Grupo] Grupo B.
        \item[Profesor] Irene Amerini.
        \item[Descripción] Examen final escrito de Computer Vision.
        \item[Fecha] junio de 2026.
        \item[Duración] 60 min.
    
    \end{description}
    \newpage

    % ------------------------------------

    \begin{observacion}
        The total sum of all the exercises is 32 points, although the maximum mark of the exam is 30.
    \end{observacion}

    \begin{ejercicio}[8 points]
    Explain the role of scale-space representation in computer vision. In your answer, discuss:
    \begin{enumerate}
        \item[a)] Why analyzing an image at multiple scales can be useful.
        \item[b)] How Gaussian filtering is used to construct a scale-space representation.
        \item[c)] The difference between smoothing at a single scale and building a Gaussian pyramid.
        \item[d)] One application where multi-scale analysis is important.
    \end{enumerate}
    \end{ejercicio}

    \begin{ejercicio}[8 points]
        Describe the Harris corner detector. In your answer, explain:
        \begin{enumerate}
            \item[a)] The intuition behind detecting corners using local intensity variations.
            \item[b)] The role of the second-moment matrix.
            \item[c)] How the corner response function is computed and interpreted.
            \item[d)] Whether the detector is robust to rotation, intensity changes, and scale changes. Justify your answer.
        \end{enumerate}
    \end{ejercicio}

    \begin{ejercicio}[8 points]
        Explain the concept of epipolar geometry between two images. In your answer, discuss:
        \begin{enumerate}
            \item[a)] The meaning of epipolar lines and epipoles.
            \item[b)] How the fundamental matrix relates corresponding points in two images.
            \item[c)] Why the fundamental matrix has rank 2.
            \item[d)] How epipolar geometry can reduce the search space in stereo matching.
        \end{enumerate}
    \end{ejercicio}

    \begin{ejercicio}[8 points]
        A surveillance camera is used to monitor a fixed outdoor area. During the day, illumination changes strongly because of shadows, clouds, and sunlight. You need to design a traditional computer vision method to detect moving objects while reducing false detections caused by lighting variations. Describe:
        \begin{enumerate}
            \item[a)] Which preprocessing steps you would apply.
            \item[b)] How you would model the background and detect foreground regions.
            \item[c)] How you would distinguish real object motion from illumination changes or shadows.
        \end{enumerate}
    \end{ejercicio}
\end{document}